% !TEX TS-program = pdflatex
% CV ciblé — Ingénieur DevOps Junior (CDI/SIVP) — DATA GUIDING, El Kram, Tunis
\documentclass[a4paper,10pt]{article}
\usepackage[utf8]{inputenc}
\usepackage[T1]{fontenc}
\usepackage[scaled=0.94]{helvet}
\renewcommand{\familydefault}{\sfdefault}
\usepackage[left=1.1cm,right=1.1cm,top=0.6cm,bottom=0.5cm]{geometry}
\usepackage{enumitem}
\usepackage{titlesec}
\usepackage{xcolor}
\usepackage[hidelinks]{hyperref}
\definecolor{navy}{RGB}{23,54,93}
\pagestyle{empty}
\setlength{\parindent}{0pt}
\setlength{\parskip}{0pt}
\linespread{0.85}
\hypersetup{pdftitle={CV - Baha Eddine Belhaj Mouldi - Ingenieur DevOps Junior - Data Guiding},pdfauthor={Baha Eddine Belhaj Mouldi},pdfsubject={DevOps, CI/CD, Docker, Kubernetes, Cloud, Automatisation}}
\titleformat{\section}{\normalsize\bfseries\color{navy}}{}{0pt}{\MakeUppercase}[{\vspace{1pt}\color{navy}\hrule height 0.6pt}]
\titlespacing*{\section}{0pt}{3pt}{1.2pt}
\setlist[itemize]{leftmargin=1.1em,label=\textbullet,itemsep=0.1pt,topsep=0.2pt,parsep=0pt}
\newcommand{\cventry}[4]{\textbf{#1} \hfill \textbf{#4}\\[0.5pt]#2{\ifx\relax#3\relax\else, #3\fi}\par\vspace{1pt}}
\begin{document}
\begin{center}
{\LARGE\bfseries\color{navy} Baha Eddine Belhaj Mouldi}\\[3pt]
{\large Ingénieur DevOps Junior --- CI/CD \textbar\ Cloud \textbar\ Conteneurisation \textbar\ Automatisation}\\[4pt]
Tunis, Tunisie \quad|\quad +216 55 128 982 \quad|\quad \href{mailto:bahaeddine.belhajmouldi@gmail.com}{bahaeddine.belhajmouldi@gmail.com}\\[1pt]
\href{https://www.linkedin.com/in/bahamouldi}{linkedin.com/in/bahamouldi} \quad|\quad \href{https://github.com/bahamouldi}{github.com/bahamouldi} \quad|\quad \href{https://bahamouldi.github.io/portfolio/}{bahamouldi.github.io/portfolio}
\end{center}
\vspace{1pt}
\section{Profil}
Jeune ingénieur en génie informatique (ENIT, mention Très Bien, cycle préparatoire IPEI El Manar), passionné par l'automatisation, le déploiement continu et la culture DevOps. Expérience concrète et déjà en production : mise en place d'un pipeline CI/CD complet (Jenkins + ArgoCD GitOps) déployant un système conteneurisé sur cluster Kubernetes, avec observabilité complète (Prometheus, Grafana, ELK). À l'aise avec l'administration Linux, Docker, les manifests Kubernetes et le scripting Python/Bash. Bases solides en cloud (AWS) et premières notions Azure/Terraform. Esprit collaboratif, grande curiosité technique et forte volonté d'apprendre aux côtés d'une équipe technique. Bon niveau en français et en anglais.
\section{Compétences recherchées par le poste}
\textbf{CI/CD} : Jenkins (pipelines multi-étapes en production), ArgoCD (GitOps), GitHub Actions, notions GitLab CI --- mise en place et maintien de pipelines\\[0.5pt]
\textbf{Cloud} : AWS (EC2, S3, VPC, IAM, RDS, CloudWatch), notions Azure, exposition GCP\\[0.5pt]
\textbf{Conteneurisation} : Docker, Docker Compose, Kubernetes (Deployment, Service, Ingress, PVC, Secrets, ConfigMap) --- déploiement continu zéro downtime\\[0.5pt]
\textbf{Linux/Unix \& Scripting} : Linux (Ubuntu, Debian), Bash, Python, PowerShell --- scripts d'automatisation de déploiement et de reporting\\[0.5pt]
\textbf{Surveillance \& fiabilité} : Prometheus, Grafana, ELK Stack (Elasticsearch, Logstash, Kibana), alerting temps réel, dashboards\\[0.5pt]
\textbf{Infrastructure as Code} : Terraform (notions), manifests Kubernetes YAML, Docker Compose\\[0.5pt]
\textbf{Collaboration \& bonnes pratiques} : travail en binôme avec les développeurs, documentation technique, partage de connaissances, méthodologie Agile
\section{Expérience professionnelle}
\cventry{Ingénieur DevOps / Sécurité --- Stagiaire PFE (mention Très Bien)}{Beehive Entreprises}{Tunisie}{02/2026 -- 06/2026}
\begin{itemize}
  \item Mise en place et maintien d'un pipeline CI/CD complet : Jenkins (8 étapes) + ArgoCD (GitOps), déploiement continu zéro downtime (maxUnavailable: 0, maxSurge: 1) sur cluster Kubernetes (5 nœuds, Calico CNI).
  \item Automatisation du déploiement : manifests Kubernetes complets (Deployment, Service, Ingress, PVC, Secrets), health/liveness probes, dimensionnement des ressources.
  \item Surveillance des systèmes et fiabilité : stack d'observabilité complète (Prometheus/Grafana --- 26 panels, ELK --- Filebeat/Logstash/Elasticsearch/Kibana), alerting temps réel.
  \item Scripts d'automatisation Python/Bash pour le déploiement, le monitoring et le reporting ; documentation technique du pipeline pour les équipes.
\end{itemize}
\cventry{Développeur Full-Stack, Freelance}{Beehive Entreprises}{Tunisie}{01/2024 -- 05/2026}
\begin{itemize}
  \item Livraison de 5+ applications avec conteneurisation Docker et déploiement CI/CD Jenkins, en collaboration directe avec les équipes de développement pour optimiser le processus de livraison.
\end{itemize}
\cventry{Stagiaire Cybersécurité --- Automatisation \& Supervision SAP}{Streamlink}{Tunisie}{07/2025 -- 08/2025}
\begin{itemize}
  \item Pipelines d'automatisation Python, dashboards ELK, alerting temps réel orchestré via N8N.
\end{itemize}
\cventry{Stagiaire Sécurité Réseau (Linux \& supervision)}{Tunisie Telecom}{Tunisie}{07/2024}
\begin{itemize}
  \item Administration Linux, analyse de trafic réseau, dashboards Elastic/Kibana, 10+ vulnérabilités identifiées et documentées.
\end{itemize}
\section{Projets clés}
\cventry{Infrastructure Kubernetes en Production --- CI/CD \& Observabilité}{Docker, Kubernetes, Jenkins, ArgoCD, Prometheus, Grafana, ELK}{}{2026}
\begin{itemize}
  \item Cluster Kubernetes complet (manifests, affinité de nœuds, health probes, dimensionnement des ressources), pipeline CI/CD Jenkins + ArgoCD GitOps, monitoring temps réel. Latence overhead $<$20ms en production.
\end{itemize}
\cventry{Architecture Microservices Conteneurisée}{TypeScript, Node.js, Docker, Kubernetes, REST API}{}{2026}
\begin{itemize}
  \item Application distribuée en microservices : services découplés, orchestration de conteneurs, structure prête pour CI/CD.
\end{itemize}
\cventry{Pipeline de Streaming Temps Réel}{Kafka, PySpark, PostgreSQL, Docker Compose}{}{2024}
\begin{itemize}
  \item Architecture événementielle (6 services Docker Compose) : Kafka producer/consumer, traitement de flux PySpark, agrégation fenêtrée.
\end{itemize}
\section{Formation}
\cventry{Diplôme d'Ingénieur en Génie Informatique --- Mention Très Bien}{École Nationale d'Ingénieurs de Tunis (ENIT)}{Tunisie}{09/2023 -- 06/2026}
\begin{itemize}
  \item Classé 65e sur 600+ au concours national d'entrée. Diplômé le 25/06/2026.
\end{itemize}
\cventry{Classes préparatoires --- Physique-Technologie}{Institut Préparatoire aux Études d'Ingénieurs, El Manar}{Tunisie}{09/2021 -- 06/2023}
\section{Certifications}
Réseaux (CCNA) --- Cisco (2025) ; Git \& GitHub --- 365 Data Science (2024) ; Fondements du Deep Learning --- NVIDIA (2025) ; Machine Learning --- NVIDIA (2024).
\section{Langues}
Français (Courant) \quad|\quad Anglais (B2) \quad|\quad Arabe (Natif)
\end{document}
